\chapter{Выполнение лабораторной работы}
\label{ch:lab13}

\section{Цель работы}

Цель работы --- закрепить навыки программирования на JavaScript при создании интерактивного веб-приложения, включающего базовые вычисления, обработку строк и массивов, работу с API, динамическое обновление DOM и подготовку проекта к развёртыванию.

\section{Постановка задания}

Согласно заданию лабораторная работа №13 включает четыре основные части:
\begin{itemize}
    \item выполнение упражнений на переменные, арифметические операции и циклы;
    \item реализацию работы с массивом сообщений и поиском по строкам;
    \item создание страницы с получением данных через API и выводом карточек;
    \item подготовку проекта к запуску на Node.js-хостинге и оформление результатов в отчёте.
\end{itemize}

\section{Ход выполнения}

\subsection{Часть 1. Баланс клубных звёзд}

Первая часть лабораторной работы оформлена как история участницы аниме-клуба. В JavaScript созданы переменные с именем участницы и её стартовым балансом звёзд. Далее с помощью цикла по дням недели моделируется изменение количества звёзд: за посещение тематических ивентов начисляются бонусы, а на ежедневную покупку мерча выполняется списание баллов. Результат по каждому дню выводится в отдельный консольный блок на странице.

\subsection{Часть 2. Чат фан-клуба}

Во второй части реализован массив сообщений, стилизованный под переписку участниц клуба о любимых аниме-героинях. Сообщения выводятся последовательно с указанием авторов, а затем по ключевому слову ``вайфу'' выполняется поиск с помощью метода `includes`. Найденные фразы отображаются отдельно, что демонстрирует обработку строк и перебор массива.

\subsection{Часть 3. Каталог аниме-героинь}

Третья часть представляет собой основную визуальную секцию сайта. На сервере Express создан собственный API-маршрут `/api/heroines`, возвращающий JSON-массив с описанием аниме-героинь. На клиенте данные запрашиваются с помощью `fetch`, после чего формируются карточки с изображением, названием, аниме, архетипом и настроением персонажа. Для каждой карточки создаётся модальное окно с расширенным описанием, цитатой и дополнительными характеристиками. Таким образом, в проекте показана связка ``серверный маршрут --- клиентский запрос --- динамический DOM''.

\subsection{Часть 4. Подготовка к публикации}

Для публикации сайт оформлен как Node.js-проект с Express-сервером. Статические ресурсы расположены в каталоге `public`, а API-маршрут и служебный маршрут `/health` реализованы в файле `server.js`. Благодаря этому приложение можно запускать локально и размещать на Replit или другом совместимом хостинге без дополнительной сборки.

\section{Результат выполнения}

В результате выполнения лабораторной работы был создан единый тематический проект ``Витрина аниме-девочек''. Интерфейс оформлен в светлой аниме-стилистике с карточками героинь, модальными окнами и отдельными блоками для демонстрации базовых заданий по JavaScript. Все части лабораторной работы объединены в одном приложении и работают согласованно.

\section{Листинг ключевых файлов}

\subsection{Серверная часть проекта}
\lstinputlisting[language=Java,caption={Файл \texttt{server.js}.}]{../../replit_one_project/server.js}

\subsection{Главная страница}
\lstinputlisting[language=HTML,caption={Файл \texttt{public/index.html}.}]{../../replit_one_project/public/index.html}

\subsection{Клиентская логика JavaScript}
\lstinputlisting[language=Java,caption={Файл \texttt{public/js/app.js}.}]{../../replit_one_project/public/js/app.js}

\subsection{Стили интерфейса}
\lstinputlisting[caption={Файл \texttt{public/css/styles.css}.}]{../../replit_one_project/public/css/styles.css}

\section{Выводы}

В ходе лабораторной работы №13 были отработаны основные возможности JavaScript: использование переменных и циклов, обработка массивов и строк, работа с `fetch`, генерация HTML-разметки средствами JavaScript и организация проекта с собственным API на Express. Тематическое оформление под аниме-девочек позволило сделать итоговый проект цельным, более выразительным и наглядным.
